\section*{Chapter \ref{ch:b1:geological-maps} --- Geological Maps}

\begin{solution}{exo:b1:geological-maps:1}
\omterm{def:b1:geological-maps:attitude}{Strike}: \omterm{def:b1:spherical-astronomy:horizontal}{azimuth} of the horizontal lines of the plane; \omterm{def:b1:geological-maps:attitude}{dip}: angle between the
plane and the horizontal, measured perpendicular to the \omterm{def:b1:geological-maps:attitude}{strike}; \omterm{def:b1:geological-maps:attitude}{dip
direction}: \omterm{def:b1:spherical-astronomy:horizontal}{azimuth} of the steepest downward line. For $135/40$, the \omterm{def:b1:geological-maps:attitude}{strike} by
the right-hand rule is $135 - 90 = \ang{045}$. Symbol: a \omterm{def:b1:geological-maps:attitude}{strike} line oriented
N\,45\,E, with a tick on its south-east side (toward \ang{135}) labelled 40.
\end{solution}

\begin{solution}{exo:b1:geological-maps:2}
$z = 500 - \tan\ang{30}\,(x\sin\ang{90} + y\cos\ang{90}) = 500 - 0.577\,x$.
At $(1000, 0)$: $500 - 577 = \qty{-77}{m}$. At $(0, 1000)$: \qty{500}{m}, the
plane being horizontal along its north--south \omterm{def:b1:geological-maps:attitude}{strike}.
\end{solution}

\begin{solution}{exo:b1:geological-maps:3}
A trace that follows a contour is a horizontal contact, at \qty{350}{m}. A
trace that runs straight across hills and valleys belongs to a vertical
contact.
\end{solution}

\begin{solution}{exo:b1:geological-maps:4}
Along the river, the core of the \omterm{def:b1:geological-maps:fold}{anticline} shows the schist and the granite
(beneath the sandstone); the core of the \omterm{def:b1:geological-maps:fold}{syncline} is occupied by the chalk.
The oldest units lie lowest in the sequence: where the layers arch upward
and erosion cuts deepest into them, the oldest appear; where they sag, the
youngest are preserved in the middle.
\end{solution}

\begin{solution}{exo:b1:geological-maps:5}
The first two points are at the same elevation: the line $y = 0$ is a \omterm{def:b1:geological-maps:attitude}{strike}
line, so the \omterm{def:b1:geological-maps:attitude}{strike} is east--west. The contact is \qty{100}{m} lower
\qty{500}{m} to the north: it \omterm{def:b1:geological-maps:attitude}{dips} north, $\tan\delta = 100/500 = 0.2$,
$\delta \approx \ang{11.3}$; \omterm{def:b1:geological-maps:attitude}{dip direction} \ang{000}, \omterm{def:b1:geological-maps:attitude}{strike} \ang{270} by the
right-hand rule. With the theorem: $\overrightarrow{P_1P_2} = (1000, 0, 0)$,
$\overrightarrow{P_1P_3} = (0, 500, -100)$, $\vec n = (0, 100\,000,
500\,000)$, $\cos\delta = 500\,000/509\,902 \approx 0.981$, $\delta \approx
\ang{11.3}$, and $(n_x, n_y) = (0, 100\,000)$ points north: the same result.
\end{solution}

\begin{solution}{exo:b1:geological-maps:6}
Downstream is north, $s = 0.03$ (\ang{1.72}). (a) $a = 0 < s$: V upstream
(south), along the contours. (b) $a = -\tan\ang{20} < 0 < s$: V upstream.
(c) $a = \tan\ang{1} \approx 0.017 < s$: V upstream. (d) $a = \tan\ang{20}
\approx 0.364 > s$: V downstream (north).
\end{solution}

\begin{solution}{exo:b1:geological-maps:7}
$\xi_{\text{floor}} - \xi_{\text{side}} = h/(a - s) = 80/(0.364 - 0.03)
\approx \qty{240}{m}$: the trace on the side is about \qty{240}{m} upstream of
the apex.
\end{solution}

\begin{solution}{exo:b1:geological-maps:8}
The \omterm{def:b1:geological-maps:fold}{anticline}'s closure points in the direction of plunge: the axis plunges
north. With an axis plunging north, a \omterm{def:b1:geological-maps:fold}{syncline}'s \omterm{def:b1:geological-maps:geological-map}{outcrops} close to the
south.
\end{solution}

\begin{solution}{exo:b1:geological-maps:9}
The east-south-east side shows younger units (chalk and shale against older
units opposite them): that block moved down. The ticks of the fault symbol
are on that side, the downthrown side.
\end{solution}

\begin{solution}{exo:b1:geological-maps:10}
A line common to both planes is perpendicular to both normals, hence
parallel to $\vec n_1 \times \vec n_2$. For $090/30$: $\vec n_1 = (0.577, 0,
1)$; for $270/30$: $\vec n_2 = (-0.577, 0, 1)$. Then $\vec n_1 \times \vec n_2
= (0\cdot 1 - 1\cdot 0,\ 1\cdot(-0.577) - 0.577\cdot 1,\ 0) = (0, -1.155,
0)$: a horizontal north--south line. The first plane \omterm{def:b1:geological-maps:attitude}{dips} east, the second
west: they \omterm{def:b1:geological-maps:attitude}{dip} away from the line, which is the crest of an \omterm{def:b1:geological-maps:fold}{anticline} with a
horizontal axis.
\end{solution}

\begin{solution}{exo:b1:geological-maps:11}
The vertical dyke's trace is the straight line through the origin trending
\ang{045}, whatever the elevation: on the hill it lies somewhere on that same
line. The dipping dyke is the plane $z = -\tan\ang{60}\,d$, where $d$ is the
horizontal distance from the \omterm{def:b1:geological-maps:attitude}{strike} line through the origin, counted
positive toward the south-east. At $z = \qty{300}{m}$: $d = -300/\tan\ang{60}
\approx \qty{-173}{m}$: the trace lies \qty{173}{m} north-west of the
\omterm{def:b1:geological-maps:attitude}{strike} line through the origin, on the \omterm{def:b1:geological-maps:attitude}{up-dip} side.
\end{solution}

\begin{solution}{exo:b1:geological-maps:12}
\[
  \frac{\partial z}{\partial x} \approx \frac{405 - 420}{100} = -0.15, \qquad
  \frac{\partial z}{\partial y} \approx \frac{412 - 420}{100} = -0.08,
\]
so the slope is $\arctan\sqrt{0.15^2 + 0.08^2} = \arctan 0.17 \approx
\ang{9.6}$. The ground
descends along $-\nabla z = (0.15, 0.08)$, whose \omterm{def:b1:spherical-astronomy:horizontal}{azimuth} is
$\arctan(0.15/0.08) \approx \ang{62}$: the cell faces east-north-east.
\end{solution}

\begin{solution}{pb:b1:geological-maps:1}
\textbf{1.} Schist, granite, sandstone, limestone, shale, chalk, alluvium.

\textbf{2.} From west to east: chalk and shale, limestone, sandstone, then
schist and granite near $x = \qty{3.5}{km}$, sandstone, limestone and shale
again, then chalk; alluvium along the whole floor.

\textbf{3.} The sandstone, with the schist and the granite beneath it along
the river: the oldest units in the core, an \omterm{def:b1:geological-maps:fold}{anticline}.

\textbf{4.} The chalk, the youngest unit: a \omterm{def:b1:geological-maps:fold}{syncline}.

\textbf{5.} The base of the sandstone is highest on the \omterm{def:b1:geological-maps:fold}{anticline}'s axis,
and only along the river has erosion cut down below it.

\textbf{6.} An unconformity: the sandstone rests on two different older
units, schist and granite, and cuts off the outline of the granite.

\textbf{7.} $435 - 261 = 174$, $432 - 428 = 4$, $512 - 586 = \qty{-74}{m}$.

\textbf{8.} $\overrightarrow{E_1E_2} = (600, 1200, -170)$,
$\overrightarrow{E_1E_3} = (1100, 400, -248)$.

\textbf{9.} The cross product is
\[
  (-229\,600,\ -38\,200,\ -1\,080\,000);
\]
oriented upward, $\vec n_E = (229\,600,\ 38\,200,\ 1\,080\,000)$.

\textbf{10.} $\|\vec n_E\| \approx 1\,104\,800$; $\cos\delta \approx 0.9776$,
$\delta \approx \ang{12.2}$.

\textbf{11.} \omterm{def:b1:spherical-astronomy:horizontal}{Azimuth} of $(229\,600, 38\,200)$: $\arctan(229\,600/38\,200)
\approx \ang{80.6}$; \omterm{def:b1:geological-maps:attitude}{attitude} $081/12$, \omterm{def:b1:geological-maps:attitude}{strike} $\approx \ang{351}$.

\textbf{12.} Yes: the symbol on the east limb has a \omterm{def:b1:geological-maps:attitude}{strike} line near
north--south and a \omterm{def:b1:geological-maps:attitude}{dip} tick to the east, labelled 12.

\textbf{13.} $524 - 669 = -145$, $421 - 459 = -38$, $319 - 209 =
\qty{110}{m}$.

\textbf{14.} $\overrightarrow{W_1W_2} = (700, 1200, 107)$ and
$\overrightarrow{W_1W_3} = (1100, -600, 255)$. Their cross product,
oriented upward, is
\[
  (370\,200,\ -60\,800,\ -1\,740\,000) \quad\longrightarrow\quad
  \vec n_W = (-370\,200,\ 60\,800,\ 1\,740\,000).
\]

\textbf{15.} $\cos\delta = 1\,740\,000/\|\vec n_W\| \approx 0.9776$, $\delta
\approx \ang{12.2}$; \omterm{def:b1:geological-maps:attitude}{dip direction}: \omterm{def:b1:spherical-astronomy:horizontal}{azimuth} of $(-370\,200, 60\,800)$,
$\approx \ang{279.3}$. \omterm{def:b1:geological-maps:attitude}{Attitude} $279/12$, \omterm{def:b1:geological-maps:attitude}{strike} $\approx \ang{189}$: the west
limb \omterm{def:b1:geological-maps:attitude}{dips} west-north-west, the east limb east-north-east.

\textbf{16.} $\vec n_E \times \vec n_W = (804 \times 10^6,\ -799\,320 \times
10^6,\ 28\,101 \times 10^6)$ (rounded).

\textbf{17.} The downward direction is $(-804, 799\,320, -28\,101) \times
10^6$: trend $\arctan(-804/799\,320) \approx \ang{0}$, that is north; plunge
$\arctan\bigl(28\,101/\sqrt{804^2 + 799\,320^2}\bigr) \approx \ang{2.0}$. The
axis plunges \ang{2} toward the north.

\textbf{18.} $8000 \times \tan\ang{2.0} \approx \qty{280}{m}$ lower, near
\qty{-50}{m}. The \omterm{def:b1:geological-maps:fold}{anticline}'s core is highest near the river and sinks to
the north, so its \omterm{def:b1:geological-maps:geological-map}{outcrop} closes to the north, as the map shows.

\textbf{19.} $\overrightarrow{F_1F_2} = (1627, 2500, -276)$ and
$\overrightarrow{F_1F_3} = (3124, 5500, 77)$ give the upward normal
\[
  (1\,710\,500,\ -987\,503,\ 1\,138\,500), \qquad
  \delta = \arccos\frac{1\,138\,500}{2\,279\,700} \approx \ang{60},
\]
and a \omterm{def:b1:geological-maps:attitude}{dip direction} $\approx \ang{120}$. \omterm{def:b1:geological-maps:attitude}{Attitude} $120/60$, \omterm{def:b1:geological-maps:attitude}{strike}
\ang{030}.

\textbf{20.} The fault \omterm{def:b1:geological-maps:attitude}{dips} toward the east-south-east, so that side is the
hanging wall, the block lying above the dipping fault plane; it moved down (younger units at the surface): a normal fault.

\textbf{21.} The fault plane \omterm{def:b1:geological-maps:attitude}{dips} east-south-east, downstream and far more
steeply than the river's floor descends: lower ground meets the plane
farther \omterm{def:b1:geological-maps:attitude}{down-dip}, so the trace swings east-south-east in the valley --- the
V-rule applied to the fault plane.

\textbf{22.} Intrusion of the granite; erosion before the sandstone;
deposition of the chalk; folding; faulting; deposition of the alluvium.

\textbf{23.} After 90 million years ago (the chalk is folded) and before
2.58 million years ago (the alluvium is not faulted); the folding came
first, the faulting after it.

\textbf{24.} The fault descends \qty{1405}{m} from F1 (\qty{405}{m}) to
\qty{-1000}{m}; at a \omterm{def:b1:geological-maps:attitude}{dip} of \ang{60} that takes a horizontal distance of
$1405/\tan\ang{60} \approx \qty{811}{m}$ \omterm{def:b1:geological-maps:attitude}{down-dip}. The well must be drilled
about \textbf{\qty[text-series-to-math=true]{810}{m} east-south-east of F1} (\omterm{def:b1:spherical-astronomy:horizontal}{azimuth} \ang{120}).
\end{solution}
